\section{Pendahuluan}
\label{sec-pendahuluan}

Perkembangan infrastruktur jaringan saat ini tidak hanya berfokus pada ketersediaan layanan (\textit{availability}), melainkan juga pada kualitas pengalaman pengguna (\textit{Quality of Experience} atau QoE). Pemantauan performa jaringan konvensional umumnya dilakukan pada sisi server atau infrastruktur inti, sehingga kurang mampu menggambarkan kualitas layanan yang dirasakan langsung oleh pengguna akhir pada lapisan akses nirkabel \cite{barakabitze2020}. Untuk mengatasi celah ini, konsep \textit{Digital Experience Monitoring} (DEM) banyak diterapkan, salah satunya melalui pemantauan aktif menggunakan perangkat sensor \textit{edge} khusus seperti HPE Aruba UXI (\textit{User Experience Insight}) yang mensimulasikan aktivitas pengguna untuk mengidentifikasi gangguan secara proaktif.

Namun, solusi komersial seperti HPE Aruba UXI memerlukan biaya pengadaan dan lisensi yang relatif tinggi, serta dioptimalkan untuk skala \textit{enterprise}. Hal ini membatasi penerapannya pada lingkungan jaringan skala kecil dan menengah seperti laboratorium atau kantor cabang kecil. Sebagai alternatif, penggunaan perangkat komputer papan tunggal (\textit{single-board computer} atau SBC) berbiaya rendah mulai diteliti, seperti UXI-Lite berbasis Raspberry Pi \cite{baswara2026}. Meskipun demikian, implementasi pemantauan pada perangkat tepi murah memiliki tantangan berupa keterbatasan sumber daya komputasi, memori, dan penyimpanan. Pemantauan periodik dengan interval rapat dapat meningkatkan utilisasi CPU dan mempercepat keausan media penyimpanan akibat penulisan log secara konstan, sedangkan interval yang terlalu longgar berisiko melewatkan gangguan singkat (\textit{transient faults}).

Untuk mengatasi hal tersebut, penelitian ini merancang konsep ``Micro-UXI'', yaitu sistem pemantauan aktif berbasis perangkat \textit{edge} berbiaya rendah (Arduino Uno Q) yang berfungsi sebagai \textit{Network Experience Black-Box} ``\textit{Event Recorder}'' \cite{kucera2020}. Sistem ini melakukan pengukuran metrik secara ringan pada kondisi normal, namun secara otomatis akan merekam berkas bukti insiden (\textit{evidence bundle}) secara mendalam ketika terdeteksi adanya penurunan performa (\textit{event-trigger}). Data performa sesaat sebelum deteksi (\textit{pre-event window}) disimpan menggunakan struktur data \textit{ring buffer} pada memori RAM \cite{marwedel2021} agar tidak membebani media penyimpanan fisik.

Tantangan utama dalam implementasi Micro-UXI adalah menjaga akurasi deteksi dengan beban komputasi yang minimal pada perangkat komputasi terbatas. Penelitian ini membandingkan dua metode deteksi, yaitu \textbf{Metode Baseline} menggunakan batas statis (\textit{static threshold}) berbasis persentil ke-99 ($P_{99}$) dan \textbf{Metode Event-Driven} dengan batas dinamis menggunakan algoritma \textit{Exponential Weighted Moving Average} (EWMA) \cite{miranda2022}. Evaluasi kedua metode dilakukan melalui eksperimen suntikan gangguan (\textit{fault injection}) terkontrol \cite{cotroneo2022} untuk enam skenario gangguan (S1--S6): \textit{DNS degraded, DNS timeout burst, loss burst, high RTT, HTTP slow,} dan \textit{connectivity flap}. Kinerja sistem dinilai berdasarkan metrik \textit{Precision, Recall, F1-Score, Mean Time to Detect} (MTTD), \textit{system overhead} (CPU, RAM, bandwidth), dan kelengkapan bukti (\textit{Reconstruction Completeness Score} atau RCS). Hasil penelitian ini diharapkan dapat menjadi referensi dalam merancang arsitektur pemantauan jaringan nirkabel yang efisien pada perangkat dengan keterbatasan sumber daya.